\section{Identificação, ajuste e comparação de modelos}

Foram ajustados modelos ARIMA e SARIMA com diferenciação regular $d=1$ e, para os modelos sazonais, diferenciação sazonal $D=1$ com período $s=12$. A forma geral considerada para os modelos sazonais pode ser escrita como

\[
\Phi(B^{12})\phi(B)(1-B)^d(1-B^{12})^D y_t =
\Theta(B^{12})\theta(B)\varepsilon_t,
\]

em que $B$ é o operador de defasagem, $\phi(B)$ e $\theta(B)$ representam os polinômios autorregressivo e de média móvel, $\Phi(B^{12})$ e $\Theta(B^{12})$ representam os componentes sazonais, e $\varepsilon_t$ é o erro do modelo.

Para avaliar risco preditivo, a amostra foi dividida em treino e teste. O treino cobre janeiro de 2000 a março de 2024, com 291 observações. O teste cobre abril de 2024 a março de 2026, com 24 observações. Os modelos foram ajustados apenas no treino e comparados por previsões para o período de teste.

\begin{table}[H]
\centering
\caption{Modelos concorrentes com melhor desempenho preditivo}
\begin{tabular}{lrrrr}
\toprule
Modelo & AIC treino & BIC treino & RMSE teste & MAPE teste \\
\midrule
SARIMA(1,1,3)(0,1,1)$_{12}$ & 1238,25 & 1259,66 & 1,78 & 1,35\% \\
SARIMA(3,1,1)(0,1,1)$_{12}$ & 1245,61 & 1267,07 & 1,79 & 1,36\% \\
SARIMA(2,1,1)(0,1,1)$_{12}$ & 1243,62 & 1261,50 & 1,79 & 1,36\% \\
SARIMA(1,1,1)(1,1,1)$_{12}$ & 1243,62 & 1261,50 & 1,79 & 1,36\% \\
SARIMA(1,1,1)(0,1,1)$_{12}$ & 1241,63 & 1255,93 & 1,79 & 1,36\% \\
\bottomrule
\end{tabular}
\end{table}

O modelo com menor RMSE no teste foi o \sarimafinal. O desempenho dos cinco melhores modelos foi muito próximo, mas a escolha pelo menor erro preditivo é coerente com o objetivo de previsão do trabalho. O AIC/BIC foi usado como critério auxiliar, não como critério exclusivo.

